\documentclass[10pt,a4paper]{amsart}
\usepackage{amsmath,amssymb,amsthm,graphicx}
\usepackage{verbatim}
\newenvironment{reference}{\comment}{\endcomment}
%\newenvironment{reference}{}{}
\newenvironment{slogan}{\comment}{\endcomment}
\newenvironment{history}{\comment}{\endcomment}

% For commutative diagrams we use Xy-pic
\usepackage[all]{xy}

% We use 2cell for 2-commutative diagrams.
\xyoption{2cell}
% Theorem environments.
%
\theoremstyle{plain}
\newtheorem{theorem}[subsection]{Theorem}
\newtheorem{proposition}[subsection]{Proposition}
\newtheorem{lemma}[subsection]{Lemma}

\theoremstyle{definition}
\newtheorem{definition}[subsection]{Definition}
\newtheorem{example}[subsection]{Example}
\newtheorem{exercise}[subsection]{Exercise}
\newtheorem{situation}[subsection]{Situation}

\theoremstyle{remark}
\newtheorem{remark}[subsection]{Remark}
\newtheorem{remarks}[subsection]{Remarks}

\numberwithin{equation}{subsection}

% Macros
%
\def\lim{\mathop{\mathrm{lim}}\nolimits}
\def\colim{\mathop{\mathrm{colim}}\nolimits}
\def\Spec{\mathop{\mathrm{Spec}}}
\def\Hom{\mathop{\mathrm{Hom}}\nolimits}
\def\Ext{\mathop{\mathrm{Ext}}\nolimits}
\def\SheafHom{\mathop{\mathcal{H}\!\mathit{om}}\nolimits}
\def\SheafExt{\mathop{\mathcal{E}\!\mathit{xt}}\nolimits}
\def\Sch{\mathit{Sch}}
\def\Mor{\mathop{\mathrm{Mor}}\nolimits}
\def\Ob{\mathop{\mathrm{Ob}}\nolimits}
\def\Sh{\mathop{\mathit{Sh}}\nolimits}
\def\NL{\mathop{N\!L}\nolimits}
\def\CH{\mathop{\mathrm{CH}}\nolimits}
\def\proetale{{pro\text{-}\acute{e}tale}}
\def\etale{{\acute{e}tale}}
\def\QCoh{\mathit{QCoh}}
\def\Ker{\mathop{\mathrm{Ker}}}
\def\Im{\mathop{\mathrm{Im}}}
\def\Coker{\mathop{\mathrm{Coker}}}
\def\Coim{\mathop{\mathrm{Coim}}}

% Boxtimes
%
\DeclareMathSymbol{\boxtimes}{\mathbin}{AMSa}{"02}

%
% Macros for moduli stacks/spaces
%
\def\QCohstack{\mathcal{QC}\!\mathit{oh}}
\def\Cohstack{\mathcal{C}\!\mathit{oh}}
\def\Spacesstack{\mathcal{S}\!\mathit{paces}}
\def\Quotfunctor{\mathrm{Quot}}
\def\Hilbfunctor{\mathrm{Hilb}}
\def\Curvesstack{\mathcal{C}\!\mathit{urves}}
\def\Polarizedstack{\mathcal{P}\!\mathit{olarized}}
\def\Complexesstack{\mathcal{C}\!\mathit{omplexes}}
% \Pic is the operator that assigns to X its picard group, usage \Pic(X)
% \Picardstack_{X/B} denotes the Picard stack of X over B
% \Picardfunctor_{X/B} denotes the Picard functor of X over B
\def\Pic{\mathop{\mathrm{Pic}}\nolimits}
\def\Picardstack{\mathcal{P}\!\mathit{ic}}
\def\Picardfunctor{\mathrm{Pic}}
\def\Deformationcategory{\mathcal{D}\!\mathit{ef}}
\usepackage{cleveref}
\crefname{originalreference}{original reference}{original references}
\Crefname{originalreference}{Original reference}{Original references}
\makeatletter
\newlabel{situation-setup}{{02QL}{source}{}{sourcepointer-1}{}}
\newlabel{situation-setup@cref}{{[originalreference][0][]02QL}{[1][0][]source}{}{}{}}
\makeatother
\crefname{originallemma}{original lemma}{original lemmas}
\Crefname{originallemma}{Original lemma}{Original lemmas}
\makeatletter
\newlabel{lemma-second-adams-operator}{{0FEJ}{source}{}{sourcepointer-2}{}}
\newlabel{lemma-second-adams-operator@cref}{{[originallemma][0][]0FEJ}{[1][0][]source}{}{}{}}
\makeatother
\crefname{originalreference}{original reference}{original references}
\Crefname{originalreference}{Original reference}{Original references}
\makeatletter
\newlabel{remark-chern-classes-K}{{0FEL}{source}{}{sourcepointer-3}{}}
\newlabel{remark-chern-classes-K@cref}{{[originalreference][0][]0FEL}{[1][0][]source}{}{}{}}
\makeatother
\crefname{originalreference}{original reference}{original references}
\Crefname{originalreference}{Original reference}{Original references}
\makeatletter
\newlabel{remark-chern-character-K}{{0FEM}{source}{}{sourcepointer-4}{}}
\newlabel{remark-chern-character-K@cref}{{[originalreference][0][]0FEM}{[1][0][]source}{}{}{}}
\makeatother
\crefname{originalreference}{original reference}{original references}
\Crefname{originalreference}{Original reference}{Original references}
\makeatletter
\newlabel{remark-extend-to-finite-locally-free}{{0ESW}{source}{}{sourcepointer-5}{}}
\newlabel{remark-extend-to-finite-locally-free@cref}{{[originalreference][0][]0ESW}{[1][0][]source}{}{}{}}
\makeatother
\crefname{originallemma}{original lemma}{original lemmas}
\Crefname{originallemma}{Original lemma}{Original lemmas}
\makeatletter
\newlabel{lemma-splitting-principle}{{02UL}{source}{}{sourcepointer-6}{}}
\newlabel{lemma-splitting-principle@cref}{{[originallemma][0][]02UL}{[1][0][]source}{}{}{}}
\makeatother
\title{The second Adams operator multiplies the Chern classes and the Chern character by powers of two}
\author{The Stacks Project Authors}
\date{stacks-project:chow.tex:local-snapshot-20261012; GNU FDL 1.2 or later}
\begin{document}
\maketitle
\noindent\textit{Editorial note.} All delivered mathematical text is the author's own, copied from the pinned Stacks Project chapter file chow.tex; no statement or proof step was rewritten, reordered or removed.\par\medskip
\noindent\textit{Editorial note.} Preamble provenance: the pinned chapter file opens with an input of the shared preamble and carries no preamble text of its own, so the project's single Stacks-wide preamble staged at sources/stacks-chow/preamble.tex is excerpted verbatim as two spans, lines 14-18 (Xy-pic setup) and lines 154-217 (the native theorem declarations and mathematical macros). Nothing was added, removed or redefined; the wrapper class and the amsmath/amssymb/amsthm/graphicx base packages are supplied by the assembler.\par\medskip
\noindent\textit{Editorial note.} The counted statement and proof contain no cross reference, no citation and no included file, so the delivered body is self contained.\par\medskip
\noindent\textit{Editorial note.} Canonical location of the counted proposition: Stacks tag 0FEN (sources/stacks-chow/pinned-original-source.tex line 12902).\par\medskip
\noindent\textit{Allowed standard prerequisites (editorial).}\begin{itemize}
\item The second Adams operator and the Chern classes and character. Same chapter, Lemma labelled lemma-second-adams-operator, cited in the statement; Remarks labelled remark-chern-classes-K and remark-chern-character-K, cited in the statement. Fixes the objects related by the lemma.
\item The splitting principle. Same chapter, Lemma labelled lemma-splitting-principle, cited at proof lines 12924-12928. Reduces the computation to invertible modules after an injective pullback on Chow rings.
\item Extension of the Chern classes from constant rank to all finite locally free modules. Same chapter, Remark labelled remark-extend-to-finite-locally-free, cited at proof lines 12920-12922. Justifies assuming E has constant rank.
\end{itemize}\medskip
\ifdefined\hypertarget\hypertarget{sourcepointer-1}{}\fi
\noindent\textit{Source reference.} Original reference 02QL: sources/stacks-chow/pinned-original-source.tex, source lines 1372--1372. Same-chapter pointer into the author environment situation at source line 1372; invoked at proof line None of the counted proof. It is retained only as a navigation pointer; every load-bearing step is inside the delivered span.\par\medskip
\ifdefined\hypertarget\hypertarget{sourcepointer-2}{}\fi
\noindent\textit{Source reference.} Original lemma 0FEJ: sources/stacks-chow/pinned-original-source.tex, source lines 12753--12753. Same-chapter pointer into the author environment lemma at source line 12753; invoked at proof line None of the counted proof. It is retained only as a navigation pointer; every load-bearing step is inside the delivered span.\par\medskip
\ifdefined\hypertarget\hypertarget{sourcepointer-3}{}\fi
\noindent\textit{Source reference.} Original reference 0FEL: sources/stacks-chow/pinned-original-source.tex, source lines 12871--12871. Same-chapter pointer into the author environment remark at source line 12871; invoked at proof line None of the counted proof. It is retained only as a navigation pointer; every load-bearing step is inside the delivered span.\par\medskip
\ifdefined\hypertarget\hypertarget{sourcepointer-4}{}\fi
\noindent\textit{Source reference.} Original reference 0FEM: sources/stacks-chow/pinned-original-source.tex, source lines 12887--12887. Same-chapter pointer into the author environment remark at source line 12887; invoked at proof line None of the counted proof. It is retained only as a navigation pointer; every load-bearing step is inside the delivered span.\par\medskip
\ifdefined\hypertarget\hypertarget{sourcepointer-5}{}\fi
\noindent\textit{Source reference.} Original reference 0ESW: sources/stacks-chow/pinned-original-source.tex, source lines 7632--7632. Same-chapter pointer into the author environment remark at source line 7632; invoked at proof line 12920 of the counted proof. It is retained only as a navigation pointer; every load-bearing step is inside the delivered span.\par\medskip
\ifdefined\hypertarget\hypertarget{sourcepointer-6}{}\fi
\noindent\textit{Source reference.} Original lemma 02UL: sources/stacks-chow/pinned-original-source.tex, source lines 8429--8429. Same-chapter pointer into the author environment lemma at source line 8429; invoked at proof line 12925 of the counted proof. It is retained only as a navigation pointer; every load-bearing step is inside the delivered span.\par\medskip
\begin{lemma}
\label{lemma-adams-and-chern}
Let $(S, \delta)$ be as in Situation \ref{situation-setup}.
Let $X$ be locally of finite type over $S$. If $\psi^2$ is
as in Lemma \ref{lemma-second-adams-operator} and $c$ and $ch$ are as in
Remarks \ref{remark-chern-classes-K} and \ref{remark-chern-character-K}
then we have $c_i(\psi^2(\alpha)) = 2^i c_i(\alpha)$ and
$ch_i(\psi^2(\alpha)) = 2^i ch_i(\alpha)$
for all $\alpha \in K_0(\textit{Vect}(X))$.
\end{lemma}

\begin{proof}
Observe that the map $\prod_{i \geq 0} A^i(X) \to \prod_{i \geq 0} A^i(X)$
multiplying by $2^i$ on $A^i(X)$ is a ring map. Hence, since $\psi^2$
is also a ring map, it suffices to prove the formulas for additive generators
of $K_0(\textit{Vect}(X))$. Thus we may assume $\alpha = [\mathcal{E}]$
for some finite locally free $\mathcal{O}_X$-module $\mathcal{E}$.
By construction of the Chern classes of $\mathcal{E}$ we immediately
reduce to the case where $\mathcal{E}$ has constant rank $r$, see
Remark \ref{remark-extend-to-finite-locally-free}.
In this case, we can choose a projective smooth morphism $p : P \to X$
such that restriction $A^*(X) \to A^*(P)$ is injective
and such that $p^*\mathcal{E}$ has a finite filtration whose
graded parts are invertible $\mathcal{O}_P$-modules $\mathcal{L}_j$, see
Lemma \ref{lemma-splitting-principle}. Then
$[p^*\mathcal{E}] = \sum [\mathcal{L}_j]$ and hence
$\psi^2([p^\mathcal{E}]) = \sum [\mathcal{L}_j^{\otimes 2}]$
by definition of $\psi^2$. Setting $x_j  = c_1(\mathcal{L}_j)$
we have
$$
c(\alpha) = \prod (1 + x_j)
\quad\text{and}\quad
c(\psi^2(\alpha)) = \prod (1 + 2 x_j)
$$
in $\prod A^i(P)$ and we have
$$
ch(\alpha) = \sum \exp(x_j)
\quad\text{and}\quad
ch(\psi^2(\alpha)) = \sum \exp(2 x_j)
$$
in $\prod A^i(P)$. From these formulas the desired result follows.
\end{proof}

\end{document}
